% Introduction to the EU – Midterm All‑in‑One (History, Institutions, Readings, Exam Tips)
% Compile with: pdflatex exam-all-in-one.tex

\documentclass[8pt,landscape,a4paper]{article}
\usepackage[utf8]{inputenc}
\usepackage[T1]{fontenc}
\usepackage[left=0.8cm,right=2.6cm,top=0.8cm,bottom=0.8cm,marginparwidth=2cm,marginparsep=0.3em]{geometry}
\usepackage{multicol}
\usepackage{marginnote}
\usepackage{enumitem}
\setlist{nosep,leftmargin=*}
\usepackage{booktabs}
\usepackage{parskip}
\usepackage{microtype}
\usepackage{xcolor}
\usepackage{fancyhdr}
\definecolor{eublue}{RGB}{0,51,153}
\definecolor{charcoal}{RGB}{51,51,51}
\definecolor{softgrey}{RGB}{240,240,240}
\definecolor{notegrey}{RGB}{90,90,90}
\newcommand{\handnote}[1]{\marginnote{\raggedright\footnotesize\color{notegrey}\textit{#1}\vspace{0.4em}}}
\newcommand{\hlpink}[1]{\textcolor{charcoal}{\textbf{#1}}}
\newcommand{\hlpeach}[1]{\textcolor{eublue}{\textbf{#1}}}
\setlength{\parskip}{0.2em}
\setlength{\columnsep}{0.5em}
\pagestyle{fancy}
\fancyhf{}
\fancyfoot[L]{\raisebox{0.8em}{\small\color{notegrey} Intro EU}}
\fancyfoot[R]{\raisebox{0.8em}{\small\color{notegrey}\thepage}}
\fancypagestyle{plain}{\fancyhf{}\fancyfoot[L]{\raisebox{0.8em}{\small\color{notegrey} Intro EU}}\fancyfoot[R]{\raisebox{0.8em}{\small\color{notegrey}\thepage}}}
\raggedright

\begin{document}
\begin{center}
\vspace{0.2em}
{\normalsize\bfseries\color{eublue} Introduction to the EU — Midterm All‑in‑One}\\[0.2em]
\footnotesize\textcolor{notegrey}{Built from your notes, readings, and slides (weeks 1–8). Use this for quick recall; in the exam you analyse texts (main idea, argument, key expression).}\\[0.35em]
\noindent\colorbox{softgrey}{\parbox{0.92\textwidth}{\centering\footnotesize\color{charcoal}
\textbf{Key:} \textbf{Definitions} \quad \textcolor{eublue}{\textbf{Frameworks}} \quad \textit{Margin notes (exam tips)}
\vspace{0.25em}}}
\end{center}
\vspace{-0.4em}

\begin{multicols}{2}

%%%%%%%%%%%%%%%%%%%%%%%%%%%%%%%%%%%%%%%%%%%%%%%%%%%%
\section*{\color{eublue}0. Exam overview \& strategy}\handnote{From Mar 16 class.}
\begin{itemize}[leftmargin=*]
\item \textbf{Not tricky, but not automatic.} Professor says exam is not a puzzle or “genius” test — but you must have studied and have an overview.
\item \textbf{Time:} About 90 minutes. Enough to read, think, plan, then write structured answers.
\item \textbf{Task:} Read short texts; identify main idea, argument, key expression; situate in the long process of integration and in the four blocks.
\item \textbf{Answer structure:} Name concept/author $\to$ define in one sentence $\to$ 2–3 specific pieces of evidence $\to$ conclude (link back to question).
\end{itemize}

\subsection*{\color{eublue}Five stages of integration}\handnote{Draw this timeline at the top of your exam sheet.}
\begin{itemize}[leftmargin=*]
\item \textbf{1945–1958} — Foundations after WWII: Marshall Plan, Council of Europe, NATO, ECSC, failed EDC.
\item \textbf{1958–1973} — \hlpink{Treaty of Rome} to \hlpink{oil crisis}. EEC + Euratom, customs union, common market, CAP, Luxembourg Compromise, first enlargement (1973).
\item \textbf{1973–1992} — Crises \& adaptation: oil shock, stagflation, Southern enlargement, \hlpink{SEA} (single market), path to \hlpink{Maastricht}.
\item \textbf{1992–present} — EU, euro, pillars, Schengen, Lisbon, big enlargements, Green Deal, migration, identity.
\item \textbf{2008 crisis} — Global financial \& euro crisis as a new phase inside this long period.
\end{itemize}

\subsection*{\color{eublue}How to use this in answers}
\begin{itemize}[leftmargin=*]
\item \textbf{Always say which stage you’re in.} “This text belongs to the 1958–1973 period (Rome to oil crisis)\dots”
\item Match each treaty / policy / reading to one of these stages and one \emph{challenge} in that stage.
\end{itemize}

%%%%%%%%%%%%%%%%%%%%%%%%%%%%%%%%%%%%%%%%%%%%%%%%%%%%
\section*{\color{eublue}1. Four blocks of the course}\handnote{From syllabus and week‑6 midterm.tex.}
\begin{itemize}[leftmargin=*]
\item \textbf{Block 1 — History \& origins.} Schuman, Rome, ECSC/EEC, crises, Southern enlargement, SEA, Maastricht, enlargements.
\item \textbf{Block 2 — Institutions.} Executive, legislative, judiciary, “twin executive” (European Council + Commission), Parliament, Council, CJEU, ECB, agencies.
\item \textbf{Block 3 — Policies.} Euro, EMU, CAP, identity, citizenship, migration, Green Deal.
\item \textbf{Block 4 — Foreign \& security.} CFSP, NATO context, external dimension.
\end{itemize}

%%%%%%%%%%%%%%%%%%%%%%%%%%%%%%%%%%%%%%%%%%%%%%%%%%%%
\section*{\color{eublue}2. History — milestones (from midterm.tex)}\handnote{Know the sequence, not every date.}
\begin{itemize}[leftmargin=*]
\item \textbf{1945–1954} — Marshall Plan, OEEC, Council of Europe (1949), NATO (1949), ECSC (1951), EDC fails (1954).
\item \textbf{Schuman Plan} (9 May 1950) — Coal \& steel under common High Authority. “Europe Day.”
\item \textbf{Treaty of Paris} (1951) — ECSC. First supranational transfer (FR, DE, IT, BE, NL, LU).
\item \textbf{Treaty of Rome} (1957) — EEC + Euratom. Customs union, common market.
\item \textbf{Empty chair crisis} (1965–66) — de Gaulle withdraws France. Opposes own resources, QMV.
\item \textbf{Luxembourg Compromise} (1966) — “Vital national interest” veto $\to$ paralysis $\sim$15 years.
\item \textbf{1967} — ECSC/EEC/Euratom merged $\to$ European Community (EC). Customs union 1968.
\item \textbf{Hague Summit} (1969) — Pompidou. Enlargement, EMU on the table.
\item \textbf{1973} — UK, Ireland, Denmark join. Norway: referendum “no”. UK vetoed twice by de Gaulle before.
\item \textbf{Oil crisis} (1973) — OPEC weapon; prices $\times$4; stagflation; no common EEC response.
\item \textbf{Southern enlargement} — Greece (1981), Spain \& Portugal (1986); democracy consolidation.
\item \textbf{SEA} (1986) — Single market; extends QMV.
\item \textbf{Maastricht} (1992) — EU, euro path, three pillars, citizenship.
\item \textbf{Schengen} — Internal borders removed; 26 states, 400M+ people.
\item \textbf{Lisbon} (2007) — Streamlines institutions; more power to Parliament.
\end{itemize}

\subsection*{Enlargements \& criteria}
\begin{itemize}[leftmargin=*]
\item \textbf{Copenhagen criteria} — Democracy, rule of law, human rights; functioning market economy; adopt \textit{acquis}.
\item \textbf{Enlargements:} 1973 (UK, IE, DK), 1981 (GR), 1986 (ES, PT), 1995 (AT, FI, SE), 2004 (10 East), 2007 (BG, RO), 2013 (HR).
\end{itemize}

%%%%%%%%%%%%%%%%%%%%%%%%%%%%%%%%%%%%%%%%%%%%%%%%%%%%
\section*{\color{eublue}3. Institutions (twin executive, week‑7/8)}\handnote{Use with week‑7 notes and eu‑institutions.tex.}
\begin{itemize}[leftmargin=*]
\item \textbf{Executive} — European Council (heads of state/government) + Commission (proposes, implements). “Twin executive.”
\item \textbf{Legislative} — Parliament (elected, co‑decision) + Council of the EU (ministers by policy).
\item \textbf{Judiciary} — Court of Justice (supremacy of EU law since 1960s/70s).
\item \textbf{Other:} ECB, Court of Auditors, agencies (EUIPO etc.).
\item \textbf{Commission} — “Engine room”; guardian of treaties; proposes legislation; one commissioner per portfolio.
\item \textbf{Parliament} — Only directly elected body; represents citizens; co‑legislates with Council.
\item \textbf{Council of the EU} — Represents governments; adopts laws with Parliament.
\item \textbf{European Council} — Sets big political direction; not a co‑legislator.
\end{itemize}

%%%%%%%%%%%%%%%%%%%%%%%%%%%%%%%%%%%%%%%%%%%%%%%%%%%%
\section*{\color{eublue}4. Policies — CAP, EMU, Green Deal}\handnote{From readings and week‑8 sessions.}
\begin{itemize}[leftmargin=*]
\item \textbf{CAP} (1962) — Common prices, subsidies, “butter mountains”; central to FR; shows depth of integration and budget politics.
\item \textbf{Four freedoms} — Goods, services, capital, people. Single market logic.
\item \textbf{EMU/Euro} — Economic and Monetary Union; single currency; ECB; debates over core/periphery, crisis management (esp.\ after 2008).
\item \textbf{Green Deal \& deindustrialisation} — Link climate policy, industrial policy, and populism (readings on Germany, etc.).
\end{itemize}

%%%%%%%%%%%%%%%%%%%%%%%%%%%%%%%%%%%%%%%%%%%%%%%%%%%%
\section*{\color{eublue}5. Key readings (very short)}\handnote{From your reading summaries.}
\begin{itemize}[leftmargin=*]
\item \textbf{Zweig} — Pre‑war Europe that vanished; nostalgia; warning about nationalism, borders, bureaucracy.
\item \textbf{Judt} — British ambivalence; Commonwealth vs Europe; messy, not inevitable integration.
\item \textbf{Krawatzek \& Pestel} — “New histories” attentive to vernacular voices; critical of official pro‑EU story.
\item \textbf{Anderson} — Democratic deficit; “ever closer union?”; role of ECJ and unelected institutions.
\item \textbf{Vinen} — “History in fragments”; many Europes; 1989 doesn’t simply unite continent.
\item \textbf{Bijsmans} — Euroscepticism as multifaceted phenomenon; hard vs soft; drivers (economic, identity, democracy).
\item \textbf{Gilbert} — Against teleological EU histories; look at losers, alternatives, dead ends.
\item \textbf{Kenny \& Pearce} — Brexit and Anglosphere; imagined community beyond Europe.
\end{itemize}

%%%%%%%%%%%%%%%%%%%%%%%%%%%%%%%%%%%%%%%%%%%%%%%%%%%%
\section*{\color{eublue}6. Last‑minute checklist}\handnote{Run this before you sleep.}
\begin{itemize}[leftmargin=*]
\item Can you draw the 5‑stage timeline and place Rome, oil crisis, SEA, Maastricht, 2008?
\item Can you explain Commission vs European Council vs Council of the EU?
\item Can you give 2–3 lines on CAP, EMU, Green Deal?
\item Can you say one sentence on each main author (Zweig, Judt, Anderson, Vinen, Bijsmans, Gilbert, Kenny \& Pearce)?
\item Do you remember the answer structure (name, define, evidence, conclude)?
\end{itemize}

%%%%%%%%%%%%%%%%%%%%%%%%%%%%%%%%%%%%%%%%%%%%%%%%%%%%
\section*{\color{eublue}7. Strategy recap (from exam‑strategy.tex)}\handnote{Everything in one place.}
\subsection*{\color{eublue}What the exam is (and isn’t)}
\begin{itemize}[leftmargin=*]
\item \textbf{Not a puzzle:} exam is \emph{not} tricky or “genius level” — but you must have studied.
\item \textbf{Task:} read short texts, identify main idea, argument, key expression; match to course themes and periods.
\item \textbf{Need:} one clear mental map of integration and institutions, not every micro‑detail.
\item \textbf{Time:} about 90 minutes — use it to read, plan, and then write structured answers.
\end{itemize}

\subsection*{\color{eublue}Timeline (5 stages) — how to use it}
\begin{itemize}[leftmargin=*]
\item Stages: 1945–1958; 1958–1973 (Rome–oil crisis); 1973–1992; 1992–present; 2008 crisis.
\item Place every agreement, institutional change, and reading into one stage.
\item When you see a text, ask: \emph{which stage? which challenge? which response?}
\item Use one “middle” event per stage as hook (Rome, oil crisis, SEA, Maastricht, 2008).
\end{itemize}

\subsection*{\color{eublue}Dimensions of integration}
\begin{itemize}[leftmargin=*]
\item \textbf{Economic} — most complete (customs union, single market, euro, CAP, competition).
\item \textbf{Political} — more powers, but strong role for states and vetoes.
\item \textbf{Military / foreign} — clear limits (NATO, sovereignty, national traditions).
\item In answers, show which dimension the text talks about and how far integration goes there.
\end{itemize}

\subsection*{\color{eublue}How to answer like the professor}
\begin{itemize}[leftmargin=*]
\item \textbf{Name the concept} first (or author / agreement).
\item \textbf{Define in one line} — short, textbook‑style phrase.
\item \textbf{Two–three pieces of evidence} — dates, quotes, examples from the text.
\item \textbf{Conclude} — link back to the question (e.g.\ “Therefore, this is a low‑cost strategy,” “Therefore, this belongs to the 1958–1973 stage,” etc.).
\item Use his technical language when relevant (heterogeneity, information asymmetry, democratic deficit, teleology, etc.).
\end{itemize}


\end{multicols}

\end{document}

